% =====================================================================
%  CARNET D'ENTRAINEMENT — FICHE 6 — THÈME 1
%  Niveau FACILE — Corrigé
% =====================================================================
\documentclass[11pt,a4paper]{article}
\usepackage[utf8]{inputenc}
\usepackage[T1]{fontenc}
\usepackage{lmodern}
\usepackage[french]{babel}
\usepackage{amsmath,amssymb,mathtools}
\usepackage{icomma}
\usepackage[version=4]{mhchem}
\usepackage{siunitx}
\sisetup{output-decimal-marker={,}}
\usepackage{xcolor}
\usepackage{colortbl}
\usepackage{tikz}
\usepackage[most]{tcolorbox}
\usepackage{enumitem}
\usepackage{array}
\usepackage{booktabs}
\usepackage{fancyhdr}
\usepackage[hidelinks]{hyperref}
\usetikzlibrary{arrows.meta,positioning,calc,shapes.geometric}
\usepackage[a4paper,margin=2.1cm,headheight=26pt,headsep=8pt]{geometry}

\definecolor{primary}{RGB}{146,95,160}
\definecolor{primarydark}{RGB}{92,52,104}
\definecolor{corcol}{RGB}{0,114,178}
\definecolor{astucecol}{RGB}{198,93,9}

\tcbset{boxbase/.style={enhanced,breakable,boxrule=0.9pt,arc=2.5pt,
   left=3mm,right=3mm,top=2.2mm,bottom=2.2mm,
   fonttitle=\bfseries,coltitle=white,toptitle=0.6mm,bottomtitle=0.6mm}}
\newtcolorbox[auto counter]{cor}[1][]{boxbase,colback=corcol!5,colframe=corcol,
   title={\textsc{Corrigé}~\thetcbcounter\ifx\relax#1\relax\relax\else\textmd{ —\itshape\ #1}\fi}}
\newtcolorbox{astucebox}[1][]{boxbase,colback=astucecol!8,colframe=astucecol,
   title={\textsc{Astuce}\ifx\relax#1\relax\relax\else\textmd{ : \itshape #1}\fi}}

\pagestyle{fancy}\fancyhf{}
\fancyhead[L]{\small\textcolor{primarydark}{\textbf{Fiche 6 · Thème 1} — Corrigé}}
\fancyhead[R]{\small\textcolor{primarydark}{\textsc{Facile}}}
\fancyfoot[C]{\small\thepage}
\renewcommand{\headrulewidth}{0.8pt}
\renewcommand{\headrule}{\hbox to\headwidth{\color{primary}\leaders\hrule height \headrulewidth\hfill}}
\setlength{\parindent}{0pt}\setlength{\parskip}{3pt}

\newcommand{\rep}{\textbf{\textcolor{corcol}{Réponse : }}}

\begin{document}
\begin{center}
{\Large\bfseries\color{primarydark}Thème 1 — Corrigé du niveau Facile}
\end{center}
\medskip

\begin{cor}[électrons de valence]
Le nombre d'électrons de valence est le \textbf{numéro de colonne} de l'élément :
\begin{center}
\renewcommand{\arraystretch}{1.2}
\begin{tabular}{c|cccccccc}
\toprule
Atome & \ce{C} & \ce{O} & \ce{Cl} & \ce{H} & \ce{P} & \ce{Na} & \ce{S} & \ce{N}\\
$N_v$ & 4 & 6 & 7 & 1 & 5 & 1 & 6 & 5\\
\bottomrule
\end{tabular}
\end{center}
\rep \ce{C}=4, \ce{O}=6, \ce{Cl}=7, \ce{H}=1, \ce{P}=5, \ce{Na}=1, \ce{S}=6, \ce{N}=5.
\end{cor}

\begin{cor}[duet ou octet ?]
Seuls l'hydrogène et l'hélium visent le \textbf{duet} (2 électrons, configuration de \ce{He}) ;
tous les autres visent l'\textbf{octet} (8 électrons).
\begin{itemize}[leftmargin=1.6em,itemsep=0pt]
  \item \textbf{Duet (2)} : \ce{H}, \ce{He}.
  \item \textbf{Octet (8)} : \ce{O}, \ce{C}, \ce{F}, \ce{N}.
\end{itemize}
\end{cor}

\begin{cor}[nommer le type de liaison]
\begin{enumerate}[label=\alph*),leftmargin=2em,itemsep=1pt]
  \item \ce{Cl2} : \textbf{covalente}. Deux atomes identiques (même électronégativité) ne peuvent
        que \emph{partager} un doublet, chacun apportant un électron.
  \item \ce{NaCl} : \textbf{ionique}. Métal (\ce{Na}) + non-métal (\ce{Cl}) : \ce{Na} cède son
        électron, on obtient \ce{Na+} et \ce{Cl-} qui s'attirent.
  \item 4\textsuperscript{e} liaison \ce{N-H} de \ce{NH4+} : \textbf{covalente de coordination
        (dative)}. L'azote fournit \emph{les deux} électrons ; le proton \ce{H+} n'en apporte aucun.
  \item \ce{H2} : \textbf{covalente}. Deux \ce{H} identiques partagent un doublet (chacun atteint son duet).
\end{enumerate}
\end{cor}

\begin{cor}[doublets liants et non liants]
Autour de l'oxygène de \ce{H2O} :
\begin{enumerate}[label=\alph*),leftmargin=2em,itemsep=1pt]
  \item \rep \textbf{2 doublets liants} (les deux liaisons \ce{O-H}) ;
  \item \rep \textbf{2 doublets non liants} (les deux paires libres dessinées au-dessus de l'oxygène).
\end{enumerate}
Total autour de \ce{O} : $2\times 2 + 2\times 2 = 8$ électrons $\Rightarrow$ octet respecté.
\end{cor}

\begin{cor}[l'octet est-il respecté ?]
\begin{enumerate}[label=\alph*),leftmargin=2em,itemsep=1pt]
  \item L'azote forme 3 liaisons ($3\times 2 = 6\,e^-$) et porte 1 doublet non liant ($2\,e^-$) :
        au total $6 + 2 = \mathbf{8}$ électrons.
  \item \rep \textbf{Oui}, l'azote s'entoure de 8 électrons : la règle de l'octet est respectée.
\end{enumerate}
\end{cor}

\begin{astucebox}[le réflexe des doublets]
Un atome « voit » les électrons d'une liaison \emph{en entier} (les 2 électrons), même s'il n'en
a apporté qu'un. Pour vérifier l'octet, on compte donc : \;
(nombre de liaisons)$\times 2$ $+$ (nombre de doublets libres)$\times 2$.
\end{astucebox}

\end{document}
